\subsection{The intrinsic five-generator closure}
\label{subsec:nr-five-generators}

Use the normalization in \cref{lem:nr-max-idempotent}.

\begin{theorem}[Five generators in the nonradial branch]
\label{thm:nonradial-five-generators}
Let $F$ be a nonzero irreducible homogeneous harmonic cubic on a
Euclidean space $V=V_{\mathrm{ess}}(F)$, with minimal regular zero set
and nonradial quadratic weight. After a constant rescaling of $F$, let
$V=B\oplus U$ be the orthogonal splitting of
\cref{thm:nonradial-spectral}. For $x=b+u$, define
\begin{equation}\label{eq:nr-five-invariants}
 r=|b|^2,\qquad s=|u|^2,\qquad
 a=\ip{b^2}b,\qquad p=\ip{bu}u,\qquad h=|bu|^2.
\end{equation}
Then
\begin{equation}\label{eq:nr-five-algebra}
 \cA_R(F)=\R[r,s,a,p,h],\qquad
 R=r+s,\qquad F=\frac{a+3p}{6}.
\end{equation}
These five polynomials are algebraically independent. In particular,
the smallest Laplacian algebra of $F$ requires exactly five ordinary
algebra generators, of degrees $2,2,3,3,4$.
\end{theorem}

\begin{proof}
By \cref{thm:nonradial-spectral,prop:nr-eikonal-core},
both $B$ and $U$ are nonzero, $B$ is eikonal, and
$B^2\subset B$, $BU\subset U$, $U^2\subset B$.
By \cref{prop:nr-core-exactness}, irreducibility makes $B$ exact,
and the symmetric operators $M_b=L_b|_U$ satisfy
\begin{equation}\label{eq:nr-five-traces}
 \tr_B L_b=0,\qquad \tr_U M_b=0\qquad(b\in B).
\end{equation}
Write $m=\dim U$, $n=\dim B=m+2$, and $w=u^2\in B$.
The identities of
\cref{prop:nr-representation-identities,lem:nr-U-cubic}, together with
\eqref{eq:nr-mixed-polarization}, give
\begin{align}
 M_{b^2}&=6M_b^2-\frac r3\Id_U,
 &108M_b^3&=9rM_b+a\Id_U,\label{eq:nr-five-representation}\\
 |w|^2&=\frac{s^2}{9},
 &w^2&=\frac s3w,\notag\\
 M_wu&=\frac s9u,
 &18u(bu)&=3bw+sb.\label{eq:nr-five-mixed}
\end{align}
The eikonal identity on $B$ gives $b(b^2)=rb$ and $|b^2|^2=r^2$.

\emph{Differential closure.}
The five bidegrees with respect to $B\oplus U$ are, respectively,
\[
 (2,0),\quad(0,2),\quad(3,0),\quad(1,2),\quad(2,2).
\]
Thus, for any of the
five generators $f$ of bidegree $(i,j)$,
\[
 \Gamma(r,f)=2if,\qquad \Gamma(s,f)=2jf.
\]
All remaining gradient products are
\begin{equation}\label{eq:nr-five-gamma}
\begin{aligned}
 \Gamma(a,a)&=9r^2,\\
 \Gamma(a,p)&=18h-rs,\\
 \Gamma(p,p)&=\frac{s^2}{9}+4h,\\
 \Gamma(a,h)&=rp+\frac{sa}{3},\\
 \Gamma(p,h)&=\frac{rp}{3}+\frac{2sp}{9}+\frac{sa}{27},\\
 \Gamma(h,h)&=\frac{r+s}{3}h+\frac{p^2}{9}+\frac{ap}{27}.
\end{aligned}
\end{equation}
To verify the table intrinsically, metric invariance and
\eqref{eq:nr-five-mixed} give the gradients, written in their
$B\oplus U$ components,
\begin{equation}\label{eq:nr-five-gradients}
 \nabla a=(3b^2,0),\qquad
 \nabla p=(w,2M_bu),\qquad
 \nabla h=\left(\frac{bw}{3}+\frac{sb}{9},\,2M_b^2u\right).
\end{equation}
Their pairings, using \eqref{eq:nr-five-representation}, yield the
first five lines of \eqref{eq:nr-five-gamma}.
For the last line, polarize the eikonal identity on $B$ at $w$ to obtain
\[
 2w(wb)+w^2b=|w|^2b+2\ip wb w.
\]
Pairing with $b$ and using
$\ip{b^2}w=6h-rs/3$ gives
\[
 |bw|^2=\frac{s^2r}{9}+p^2-sh.
\]
Consequently
\[
 |\nabla_Bh|^2=\frac{sh}{3}+\frac{p^2}{9},\qquad
 |\nabla_Uh|^2=4\ip{M_b^4u}u
             =\frac{rh}{3}+\frac{ap}{27},
\]
which proves the final pairing.

The Laplacians are
\begin{equation}\label{eq:nr-five-laplacians}
 \Delta r=2(m+2),\qquad \Delta s=2m,\qquad
 \Delta a=\Delta p=0,\qquad
 \Delta h=\frac{mr+(m+2)s}{9}.
\end{equation}
To verify these formulas, exactness of $B$ implies
$\sum_{\alpha=1}^n e_\alpha^2=0$ for an orthonormal basis of $B$.
Summing the quadratic identity in \eqref{eq:nr-five-representation},
and separately taking its trace, gives
\[
 \sum_{\alpha=1}^n M_{e_\alpha}^2=\frac n{18}\Id_U,
 \qquad \tr M_b^2=\frac m{18}r.
\]
Consequently,
\[
 \Delta_Bh=2\sum_{\alpha=1}^n|M_{e_\alpha}u|^2=\frac n9s,
 \qquad \Delta_Uh=2\tr M_b^2=\frac m9r.
\]
The other formulas in \eqref{eq:nr-five-laplacians} follow from
\eqref{eq:nr-five-traces}. The differential chain rule now proves
that $\R[r,s,a,p,h]$ is a Laplacian algebra.

\emph{Recovery from $R$ and $F$.}
The grading gives $F=(a+3p)/6$. Applying
\eqref{eq:nr-five-gamma} and \eqref{eq:nr-five-laplacians} gives
\begin{equation}\label{eq:nr-five-gradient-square}
 S=\frac14(r-s/3)^2+4h,\qquad
 \Delta S=\frac{(10m+36)r+(2m+4)s}{9}.
\end{equation}
Since $R=r+s$ and $m>0$, we recover the two radii directly:
\begin{equation}\label{eq:nr-five-radius-recovery}
 r=\frac{9\Delta S-(2m+4)R}{8(m+4)},\qquad s=R-r.
\end{equation}
Both belong to $\cA_R(F)$, because $S=\Gamma(F,F)$ and
$\Delta S$ do. Setting $D=\Gamma(r,F)=a+p$ then gives
\[
 p=3F-\frac D2,\qquad
 a=\frac{3D}{2}-3F,\qquad
 h=\frac{\Gamma(a,p)+rs}{18}.
\]
Thus all five generators belong to $\cA_R(F)$. Since their algebra
is Laplacian and contains $R,F$, equality follows.

\emph{Algebraic independence.}
At the maximal unit idempotent $c$,
\cref{prop:nr-eikonal-core} supplies a nonzero $(-1)$-Peirce space
$Z\subset B$. Choose a unit vector $z\in Z$. Then $z^2=-c$, so at
$b=z$ one has $r=1$, $a=0$, and
\[
 \nabla_B r=2z,\qquad \nabla_B a=-3c.
\]
In particular, the differential of $(r,a)$ has rank two.
By \eqref{eq:nr-five-representation}, every eigenvalue of $M_z$ lies in
\[
 \{\lambda_1,\lambda_2,\lambda_3\}
   =\left\{-\frac1{2\sqrt3},0,\frac1{2\sqrt3}\right\}.
\]
Let their multiplicities be $k_1,k_2,k_3$, allowing zero provisionally.
The identities $\tr M_z=0$, $\tr M_z^2=m/18$, and
$k_1+k_2+k_3=m$ imply $k_1=k_2=k_3=m/3>0$.
Choose $u=u_1+u_2+u_3$ with a nonzero component in each eigenspace.
For fixed $b=z$, the map from $t_i=|u_i|^2$ to the remaining invariants
is
\[
 \begin{pmatrix}s\\p\\h\end{pmatrix}
 =
 \begin{pmatrix}
  1&1&1\\
  \lambda_1&\lambda_2&\lambda_3\\
  \lambda_1^2&\lambda_2^2&\lambda_3^2
 \end{pmatrix}
 \begin{pmatrix}t_1\\t_2\\t_3\end{pmatrix}.
\]
The Vandermonde matrix is invertible, and the differentials $dt_i$
are independent because $u_i\ne0$. Thus $(s,p,h)$ has rank three
in the $U$ directions, whereas $(r,a)$ depends only on $b$ and has
rank two. The full five-invariant map has rank five at this point.
Its components are therefore algebraically independent. Their
polynomial algebra has transcendence degree five and cannot be
generated by fewer than five elements.
\end{proof}

\begin{corollary}[Cartan reduction and ambient dimensions]
\label{cor:nr-cartan-dimensions}
Under the hypotheses and normalization of
\cref{thm:nonradial-five-generators}, the cubic
\begin{equation}\label{eq:nr-cartan-reduction}
 H=a-9p=6\bigl(\Gamma(r,F)-5F\bigr)
\end{equation}
satisfies
\[
 \Delta H=0,\qquad |\nabla H|^2=9R^2.
\]
It is orthogonally equivalent to a normalized Cartan cubic, and
\[
 \dim V\in\{8,14,26\}.
\]
Writing $C=H/6$ and $C_B=C\circ\operatorname{pr}_B$, one has
\begin{equation}\label{eq:nr-cartan-reconstruction}
 F=\frac{4C_B-C}{3}.
\end{equation}
\end{corollary}

\begin{proof}
The formula for $H$ follows from $6F=a+3p$ and
$\Gamma(r,F)=a+p$. By \eqref{eq:nr-five-gradients},
\[
 \nabla H=(3b^2-9u^2,-18bu).
\]
The identities $|b^2|^2=r^2$, $|u^2|^2=s^2/9$, and
$\ip{b^2}{u^2}=6h-rs/3$ therefore give
\[
 |\nabla H|^2
 =9r^2+9s^2-54(6h-rs/3)+324h=9(r+s)^2.
\]
Also $\Delta H=0$ by \eqref{eq:nr-five-laplacians}.

Put $m=\dim U$. The multiplicity calculation in the proof of
\cref{thm:nonradial-five-generators} gives $m/3\in\mathbb N$
and $m>0$. Since $\dim B=m+2$, it follows that
$\dim V=2m+2\ge8$.
Cartan's classification of harmonic cubic eikonals
\cite{Cartan39}, or \cite[Theorem~1.1]{TkachevCartan10},
now identifies $H$ with a Cartan cubic in dimension $5,8,14$, or
$26$. In the latter theorem, the additional reducible possibility
$x_N^3-3x_N\sum_{i<N}x_i^2$ has Laplacian $6(2-N)x_N$,
so harmonicity and $N\ge8$ exclude it. The same dimension bound
excludes $N=5$. Finally, $C_B=a/6$ and $C=(a-9p)/6$ give
\eqref{eq:nr-cartan-reconstruction}.
\end{proof}

The reduction recovers $F$ from a Cartan cubic and a compatible
splitting $B\oplus U$. Classifying $F$ up to equivalence therefore
also requires classifying these splittings.

The models in \cref{sec:nonradial-models} realize the sharp bound
and express \(r,s\) and the other invariants in coordinates.
The programs check their closure, recovery, and rank-five minor;
the universal spectral and trace arguments are proved above and in
\cref{app:cubic-structure}.
